\chapter{Docker Compose}

Одиночный контейнер удобен для экспериментов, но реальное приложение
состоит из нескольких сервисов: веб-сервер, API, база, кэш. Запускать
их длинными командами \cmd{docker run} неудобно и опасно: легко забыть
том, сеть или переменную. Docker Compose описывает всё приложение
в одном YAML-файле и управляет им одной командой.

\section{Зачем нужен Compose}

\begin{itemize}
  \item \textbf{Декларативность.} Вы описываете желаемое состояние:
        какие сервисы, образы, порты, тома и сети нужны.
  \item \textbf{Воспроизводимость.} Файл лежит в репозитории рядом
        с кодом; коллеге достаточно выполнить \cmd{docker compose up}.
  \item \textbf{Единая сеть и имена.} Compose создаёт сеть проекта,
        в которой сервисы доступны по именам.
  \item \textbf{Управление жизненным циклом.} Запуск, остановка,
        пересборка, просмотр логов и выполнение команд~--- без
        запоминания идентификаторов контейнеров.
\end{itemize}

\begin{defbox}
\textbf{Проект Compose}~--- набор контейнеров, сетей и томов,
объединённых общим именем (по умолчанию~--- именем каталога).
Все ресурсы проекта помечаются метками, поэтому Compose не трогает
чужие контейнеры.
\end{defbox}

\section{Установка и проверка}

В современных версиях Docker Compose~--- это плагин CLI, а не отдельный
бинарник \cmd{docker-compose}.

\begin{code}[language=bash]
$ docker compose version
Docker Compose version v2.32.1
\end{code}

Если команда не найдена, установите пакет
\cmd{docker-compose-plugin} (раздел~3.2) или отдельный плагин Compose.
Старые скрипты с \cmd{docker-compose} (через дефис) лучше переписать
на \cmd{docker compose}.

\section{Первый файл}

Создадим каталог проекта и файл \file{compose.yaml}.

\begin{code}[language=composeyaml]
services:
  web:
    image: nginx:1.27-alpine
    ports:
      - "8080:80"

  cache:
    image: redis:7-alpine
\end{code}

\begin{code}[language=bash]
$ docker compose up -d
[+] Running 3/3
 Network demo_default    Created
 Container demo-cache-1  Started
 Container demo-web-1    Started
$ docker compose ps
NAME           IMAGE               COMMAND                  SERVICE   STATUS    PORTS
demo-cache-1   redis:7-alpine      "docker-entrypoint.s…"   cache     Up
demo-web-1     nginx:1.27-alpine   "/docker-entrypoint.…"   web       Up        0.0.0.0:8080->80/tcp
$ curl -I http://localhost:8080
HTTP/1.1 200 OK
\end{code}

Остановка и удаление:

\begin{code}[language=bash]
$ docker compose stop          # остановить, сохранить контейнеры
$ docker compose start         # запустить снова
$ docker compose down          # удалить контейнеры и сеть
$ docker compose down -v       # ещё и тома проекта
\end{code}

\begin{warnbox}
\cmd{docker compose down -v} удаляет именованные тома проекта вместе
с данными. Это удобно для сброса стенда, но катастрофично, если
в томе лежит база, которую вы не забэкапили. Перед выполнением
проверьте \cmd{docker compose config --volumes}.
\end{warnbox}

\section{Структура файла}

Файл Compose состоит из нескольких разделов верхнего уровня.

\begin{code}[language=composeyaml]
name: demo                # имя проекта

services:                 # контейнеры
  web:
    image: nginx:1.27-alpine
    ports:
      - "8080:80"

volumes:                  # именованные тома
  pgdata:

networks:                 # сети
  appnet:
    driver: bridge
\end{code}

\begin{itemize}
  \item \cmd{name} задаёт имя проекта; без него используется имя
        каталога.
  \item \cmd{services}~--- обязательный раздел: описание контейнеров.
  \item \cmd{volumes} и \cmd{networks}~--- объявления ресурсов,
        на которые ссылаются сервисы.
  \item Ключ \cmd{version} в Compose~v2 не нужен и считается устаревшим.
\end{itemize}

\section{Описание сервиса}

\begin{code}[language=composeyaml]
services:
  api:
    image: shop-api:1.4.2          # или build
    build:
      context: ./api
      dockerfile: Dockerfile
      args:
        APP_VERSION: "1.4.2"
    container_name: shop-api       # необязательно, мешает масштабированию
    command: ["python", "app.py", "--port", "8000"]
    environment:
      LOG_LEVEL: info
      DATABASE_URL: postgres://app:${DB_PASSWORD}@db:5432/app
    env_file:
      - ./api/.env
    ports:
      - "127.0.0.1:8000:8000"
    volumes:
      - ./uploads:/app/uploads
    networks:
      - appnet
    depends_on:
      db:
        condition: service_healthy
    restart: unless-stopped
    healthcheck:
      test: ["CMD", "curl", "-fsS", "http://127.0.0.1:8000/health"]
      interval: 10s
      timeout: 3s
      retries: 5
      start_period: 10s
    profiles: ["core"]
\end{code}

\subsection{image и build}

Если указан \cmd{build}, Compose собирает образ локально. Если ещё
и \cmd{image}, собранный образ получает это имя. При
\cmd{docker compose up --build} образы пересобираются; без флага
используется существующий образ, если он есть.

\subsection{ports и expose}

\cmd{ports} публикует порты на хосте, \cmd{expose} делает их
доступными только внутри сети проекта. Синтаксис \cmd{ports} совпадает
с \flag{-p}:

\begin{code}[language=composeyaml]
    ports:
      - "8080:80"
      - "127.0.0.1:8443:443"
      - "5353:53/udp"
\end{code}

\subsection{environment, env\_file и подстановки}

\cmd{environment} задаёт переменные напрямую, \cmd{env\_file} читает
их из файла. В значениях работает подстановка переменных окружения
хоста и файла \file{.env}.

\begin{center}
\small
\begin{tabularx}{\textwidth}{@{}lX@{}}
\toprule
\textbf{Синтаксис} & \textbf{Значение} \\
\midrule
\cmd{\$\{VAR\}} & значение переменной; пустая строка, если не задана \\
\cmd{\$\{VAR:-default\}} & значение или значение по умолчанию \\
\cmd{\$\{VAR:?message\}} & ошибка с сообщением, если переменная не задана \\
\cmd{\$\$} & литеральный символ \cmd{\$} \\
\bottomrule
\end{tabularx}
\end{center}

\begin{code}
# .env
DB_PASSWORD=super-secret
APP_VERSION=1.4.2
\end{code}

\begin{code}[language=composeyaml]
    environment:
      POSTGRES_PASSWORD: ${DB_PASSWORD:?DB_PASSWORD is required}
\end{code}

\begin{tipbox}
Проверяйте итоговую конфигурацию командой
\cmd{docker compose config}: она подставляет переменные, раскрывает
файлы и показывает, что именно будет запущено. Секреты при этом
попадут в вывод~--- не публикуйте его в CI-логах.
\end{tipbox}

\subsection{depends\_on и здоровье сервисов}

Короткая форма \cmd{depends\_on: [db]} задаёт только порядок запуска,
но не ждёт готовности базы. Длинная форма с условием ждёт успешной
проверки здоровья.

\begin{code}[language=composeyaml]
    depends_on:
      db:
        condition: service_healthy
      cache:
        condition: service_started
\end{code}

\begin{warnbox}
Даже \cmd{condition: service\_healthy} не гарантирует, что база
готова принимать конкретную схему. Приложение должно уметь
переподключаться и повторять запросы. Healthcheck~--- вспомогательный
сигнал, а не замена устойчивости кода.
\end{warnbox}

\subsection{restart, resources и прочее}

\begin{code}[language=composeyaml]
    restart: unless-stopped
    deploy:
      resources:
        limits:
          cpus: "1.0"
          memory: 512M
\end{code}

Раздел \cmd{deploy} в Compose применяется частично: ограничения
ресурсов работают, а параметры вроде \cmd{replicas} имеют смысл
только в Swarm. Для обычного \cmd{docker compose up} используйте
\cmd{cpus} и \cmd{memory} в \cmd{deploy.resources.limits}~---
они поддерживаются.

\section{Полный пример}

Соберём приложение из API и PostgreSQL с томом, проверкой здоровья
и пользовательской сетью.

\begin{figure}[htbp]
\centering
\begin{tikzpicture}[
  font=\sffamily\scriptsize,
  svc/.style={draw=accent!55!black, fill=accent!7, rounded corners=2pt,
    minimum width=2.7cm, minimum height=1.0cm, align=center},
  ext/.style={draw=warm!65!black, fill=warm!8, rounded corners=2pt,
    minimum width=2.2cm, minimum height=0.9cm, align=center},
  vol/.style={draw=good!60!black, fill=good!8, cylinder, shape border rotate=90,
    aspect=0.25, minimum height=1.0cm, minimum width=2.3cm, align=center},
  arr/.style={-Latex, thick, accent!70},
]
  \node[ext, minimum width=1.9cm] (client) at (-1.4,0) {браузер};
  \node[svc, minimum width=2.6cm] (web) at (2.0,0) {web\\nginx:80};
  \node[svc, minimum width=2.6cm] (api) at (6.3,0) {api\\:8000};
  \node[svc, minimum width=2.6cm] (db) at (10.6,0) {db\\postgres:5432};
  \node[vol] (pg) at (10.6,-1.9) {pgdata};

  \draw[arr] (client) -- node[above, text=muted] {8080} (web);
  \draw[arr] (web) -- node[above, text=muted] {api:8000} (api);
  \draw[arr] (api) -- node[above, text=muted] {db:5432} (api -| db.west);
  \draw[arr] (db) -- (pg);

  \draw[dashed, accent2!70, rounded corners=6pt] (0.5,-2.7) rectangle (13.1,1.0);
  \node[anchor=south west, text=accent2!80!black] at (0.55,1.02) {сеть appnet};
\end{tikzpicture}
\caption{Приложение Compose: браузер обращается к web, web~--- к api
по имени сервиса, api~--- к базе; данные базы хранятся в томе}
\label{fig:compose}
\end{figure}

\begin{code}[language=composeyaml]
name: shop

services:
  api:
    build: ./api
    image: shop-api:1.4.2
    environment:
      DATABASE_URL: postgres://app:${DB_PASSWORD}@db:5432/app
      LOG_LEVEL: info
    ports:
      - "127.0.0.1:8000:8000"
    depends_on:
      db:
        condition: service_healthy
    healthcheck:
      test: ["CMD", "curl", "-fsS", "http://127.0.0.1:8000/health"]
      interval: 10s
      timeout: 3s
      retries: 5
      start_period: 10s
    restart: unless-stopped
    networks: [appnet]

  db:
    image: postgres:17
    environment:
      POSTGRES_DB: app
      POSTGRES_USER: app
      POSTGRES_PASSWORD: ${DB_PASSWORD:?DB_PASSWORD is required}
    volumes:
      - pgdata:/var/lib/postgresql/data
    healthcheck:
      test: ["CMD-SHELL", "pg_isready -U app -d app"]
      interval: 10s
      timeout: 5s
      retries: 5
    restart: unless-stopped
    networks: [appnet]

volumes:
  pgdata:

networks:
  appnet:
    driver: bridge
\end{code}

Запуск:

\begin{code}[language=bash]
$ export DB_PASSWORD=secret
$ docker compose up -d --build
$ docker compose ps
$ docker compose logs -f api
$ docker compose exec db psql -U app -d app -c 'SELECT 1;'
$ docker compose down
\end{code}

\begin{notebox}
Пароль через переменную окружения~--- компромисс для учебного примера.
В production используйте менеджер секретов (Vault, AWS Secrets Manager,
Docker secrets в Swarm, внешние CSI-драйверы в Kubernetes) и не храните
\file{.env} с паролями в репозитории.
\end{notebox}

\section{Команды Compose}

\begin{center}
\small
\begin{tabularx}{\textwidth}{@{}lX@{}}
\toprule
\textbf{Команда} & \textbf{Назначение} \\
\midrule
\cmd{up -d} & создать и запустить сервисы в фоне \\
\cmd{up --build} & пересобрать образы перед запуском \\
\cmd{ps} & состояние контейнеров проекта \\
\cmd{logs -f SERVICE} & логи сервиса в реальном времени \\
\cmd{exec SERVICE CMD} & выполнить команду в работающем контейнере \\
\cmd{run --rm SERVICE CMD} & разовый контейнер (миграции, тесты) \\
\cmd{stop} / \cmd{start} / \cmd{restart} & управление без удаления \\
\cmd{down} & удалить контейнеры и сеть проекта \\
\cmd{down -v} & дополнительно удалить тома \\
\cmd{config} & проверить и показать итоговый YAML \\
\cmd{build} & собрать образы \\
\cmd{pull} & скачать образы \\
\cmd{push} & отправить образы в реестр \\
\cmd{watch} & синхронизировать файлы и пересобирать \\
\bottomrule
\end{tabularx}
\end{center}

\section{Проекты, профили и переопределения}

\subsection{Имя проекта и несколько файлов}

Имя проекта задаётся ключом \cmd{name}, флагом \flag{-p} или
переменной \envvar{COMPOSE\_PROJECT\_NAME}. Оно влияет на имена
контейнеров, сетей и томов.

\begin{code}[language=bash]
$ docker compose -p staging up -d
$ docker compose -f compose.yaml -f compose.prod.yaml up -d
\end{code}

Файл \file{compose.override.yaml} подхватывается автоматически. Это
удобно для локальной разработки: основной файл описывает production,
а override добавляет привязку каталога с кодом и отладочные порты.

\begin{code}[language=composeyaml]
# compose.override.yaml
services:
  api:
    volumes:
      - ./api:/app
    environment:
      LOG_LEVEL: debug
\end{code}

\subsection{Профили}

Профили позволяют держать в одном файле необязательные сервисы:
админку, отладчик, тестовые зависимости.

\begin{code}[language=composeyaml]
services:
  adminer:
    image: adminer:4
    profiles: ["tools"]
    ports:
      - "127.0.0.1:8081:8080"
\end{code}

\begin{code}[language=bash]
$ docker compose up -d                    # adminer не запустится
$ docker compose --profile tools up -d    # запустится
$ docker compose --profile tools ps
\end{code}

\subsection{Режим разработки}

Compose умеет следить за файлами и синхронизировать их в контейнер
или пересобирать образ.

\begin{code}[language=composeyaml]
services:
  api:
    build: ./api
    develop:
      watch:
        - action: sync
          path: ./api/src
          target: /app/src
        - action: rebuild
          path: ./api/requirements.txt
\end{code}

\begin{code}[language=bash]
$ docker compose watch
\end{code}

\begin{tipbox}
Для горячей перезагрузки Node.js и Python обычно достаточно привязки
каталога с кодом (\cmd{./api:/app}) и переменной окружения, включающей
режим разработки. \cmd{watch} удобнее тем, что не требует правки
compose-файла под каждую ОС и умеет пересобирать образ при изменении
зависимостей.
\end{tipbox}

\section{Когда Compose уже мало}

Compose рассчитан на один хост. Он не умеет:

\begin{itemize}
  \item запускать контейнеры на нескольких машинах;
  \item автоматически заменять упавший узел;
  \item балансировать нагрузку между экземплярами сервиса;
  \item выполнять постепенные обновления без простоя на кластере.
\end{itemize}

Когда приложение перерастает один сервер, на смену приходят
оркестраторы: Docker Swarm (проще, встроен в Docker) или Kubernetes
(стандарт де-факто). Знания Compose при этом не пропадают: структура
сервисов, переменные, тома и сети переносятся почти один в один.

\begin{summarybox}
\begin{itemize}
  \item Compose описывает многоконтейнерное приложение в
        \file{compose.yaml}.
  \item Проект объединяет сервисы, сети и тома под общим именем.
  \item В сети проекта сервисы доступны по именам; порты публикуются
        только там, где это нужно.
  \item \cmd{depends\_on} с \cmd{service\_healthy} ждёт готовности,
        но приложение всё равно должно быть устойчивым к переподключению.
  \item Профили и override-файлы позволяют держать в одном проекте
        разные сценарии запуска.
\end{itemize}
\end{summarybox}

\section*{Вопросы для самопроверки}
\begin{selfcheck}
  \item Чем \cmd{docker compose up} отличается от набора \cmd{docker run}?
  \item Как в Compose обратиться к базе из другого сервиса?
  \item Чем \cmd{ports} отличается от \cmd{expose}?
  \item Зачем нужна длинная форма \cmd{depends\_on}?
  \item Что удаляет \cmd{docker compose down -v}?
  \item Как запустить сервис, помеченный профилем \cmd{tools}?
\end{selfcheck}
